\textbf{Hướng phát triển tương lai}

Trong các phiên bản tiếp theo, hệ thống có thể được phát triển theo các hướng sau:
\begin{itemize}
    \item Hiệu chuẩn cảm biến TDS, pH và CO2 bằng mẫu chuẩn để tăng độ chính xác của dữ liệu.
    \item Bổ sung cơ chế lưu đệm tại gateway khi mất Internet và tự động gửi bù dữ liệu sau khi kết nối phục hồi.
    \item Mở rộng hệ thống nhiều node cho nhiều phòng hoặc nhiều khu vực cây trồng trong căn hộ.
    \item Bổ sung thuật toán điều khiển tự động dựa trên ngưỡng hoặc luật điều khiển, ví dụ tự bật quạt khi CO2 cao hoặc tự bật đèn khi ánh sáng thấp.
    \item Bổ sung cảm biến VOC, bụi mịn và lưu lượng thông gió để đánh giá IAQ đầy đủ hơn.
    \item Thử nghiệm với nhiều loài cây trong nhà hoặc mô hình thủy canh nhỏ để so sánh khả năng duy trì điều kiện sinh trưởng và hỗ trợ cải thiện môi trường.
    \item Cải thiện điều khiển điều hòa bằng cách lưu trạng thái thiết bị, học mã IR của nhiều hãng điều hòa và bổ sung phản hồi nhiệt độ sau điều khiển.
    \item Nâng cấp ứng dụng Flutter với thông báo đẩy, biểu đồ lịch sử chi tiết và cấu hình ngưỡng riêng cho từng loại cây.
\end{itemize}

Nhìn chung, đề tài đã chứng minh tính khả thi của việc kết hợp cây trồng trong nhà với hệ thống IoT để giám sát và kiểm soát chất lượng không khí trong căn hộ. Hướng tiếp cận này có thể tiếp tục mở rộng thành giải pháp nhà thông minh hướng tới sức khỏe, tiết kiệm năng lượng, chăm sóc cây trồng tự động và quản lý môi trường sống bền vững.
